\section{Operadores lineales y adjuntos}
\label{OperatorsAdjoints}

La transformada de Radon y la retroproyección filtrada se describen mediante operadores lineales. En esta sección se recuerdan únicamente las nociones generales que serán necesarias más adelante; las propiedades específicas del operador de Radon se desarrollarán en el capítulo siguiente.\\

Sea $X$ un espacio vectorial normado y sea $Y$ otro espacio vectorial normado, ambos sobre $\mathbb{R}$ o $\mathbb{C}$. Una aplicación $T:X\to Y$ se llama operador lineal si
\[
T(\alpha x+\beta y)=\alpha T x+\beta T y
\]
para todo $x,y\in X$ y escalares $\alpha,\beta$. Si además existe una constante $C\geq 0$ tal que
\[
\|Tx\|_{Y}\leq C\|x\|_{X},
\qquad x\in X,
\]
se dice que $T$ es acotado. En espacios normados, para operadores lineales, acotación y continuidad son equivalentes; este hecho básico puede consultarse en Rynne y Youngson \cite[Sec.~4.1]{rynne_youngson2000}.\\

En los espacios $L^{2}$ se usa el producto interno
\[
\langle u,v\rangle_{L^{2}(\Omega)}
=
\int_{\Omega} u(x)\overline{v(x)}\,dx,
\]
con la convención análoga en dominios de una variable. Si $H_{1}$ y $H_{2}$ son espacios de Hilbert y $T:H_{1}\to H_{2}$ es lineal y acotado, su adjunto es el operador $T^{*}:H_{2}\to H_{1}$ caracterizado por
\[
\langle Tx,y\rangle_{H_{2}}
=
\langle x,T^{*}y\rangle_{H_{1}},
\qquad x\in H_{1},\; y\in H_{2}.
\]
La existencia y unicidad del adjunto se deducen del teorema de representación de Riesz; aquí se usarán como resultado estándar \cite[Sec.~5.1]{rynne_youngson2000}.\\

En tomografía, la identidad anterior se aplicará al operador de Radon una vez fijadas sus variables y medidas. En ese contexto, el operador que aparece como adjunto formal será la retroproyección. La fórmula explícita no se introduce aquí para evitar duplicación: se derivará después de definir $\mathcal{R}$ y de especificar el intervalo angular $[0,\pi)$.\\

El multiplicador asociado al filtro rampa no será acotado en todo $L^{2}$, porque su símbolo crece sin cota. Por esta razón, en la sección siguiente se estudiará como un operador no acotado definido sobre su dominio natural denso.
